
\documentclass{article}
\usepackage{colortbl}
\usepackage{makecell}
\usepackage{multirow}
\usepackage{supertabular}

\begin{document}

\newcounter{utterance}

\centering \large Interaction Transcript for game `cladder', experiment `full\_v1.5\_default', episode 2060 with Qwen3.5{-}4B.
\vspace{24pt}

{ \footnotesize  \setcounter{utterance}{1}
\setlength{\tabcolsep}{0pt}
\begin{supertabular}{c@{$\;$}|p{.15\linewidth}@{}p{.15\linewidth}p{.15\linewidth}p{.15\linewidth}p{.15\linewidth}p{.15\linewidth}}
    \# & \multicolumn{2}{c}{Player} && \multicolumn{2}{c}{Game Master} \\
    \hline

    \theutterance \stepcounter{utterance}  
    & & & \multicolumn{4}{p{0.6\linewidth}}{
        \cellcolor[rgb]{0.9,0.9,0.9}{
            \makecell[{{p{\linewidth}}}]{
                \texttt{\tiny{[P1$\langle$GM]}}
                \texttt{You are an expert in causal inference. The following question is not a typical commonsense query, but rather a meticulously designed question created by a professor specializing in causal inference, intended to assess the students' mastery of the course content.} \\
\texttt{Imagine a self{-}contained, hypothetical world with only the following conditions, and without any unmentioned factors or causal relationships: Encouragement level has a direct effect on studying habit and exam score. Studying habit has a direct effect on exam score. The overall probability of encouragement is 91\%. For students who are not encouraged, the probability of high exam score is 72\%. For students who are encouraged, the probability of high exam score is 52\%. Is high exam score more likely than low exam score overall?} \\
\texttt{Start your answer with 'Yes' or 'No', followed by additional reasoning or evidence to support your explanation.} \\
            }
        }
    }
     \\ \\

    \theutterance \stepcounter{utterance}  
    & \multicolumn{4}{p{0.6\linewidth}}{
        \cellcolor[rgb]{0.561,0.737,0.561}{
            \makecell[{{p{\linewidth}}}]{
                \texttt{\tiny{[P1$\rangle$GM]}}
                \texttt{Yes} \\
\\ 
\texttt{**Reasoning:**} \\
\\ 
\texttt{To determine whether a high exam score is more likely than a low exam score overall, we must calculate the marginal probability of a high exam score (\$P(\textbackslash{}text\{High Score\})\$) based on the provided causal structure and probabilities.} \\
\\ 
\texttt{1.  **Define the Variables and Probabilities:**} \\
\texttt{    *   Let \$E\$ be the event of being Encouraged.} \\
\texttt{    *   Let \$S\$ be the event of having a High Exam Score.} \\
\texttt{    *   \$P(E) = 0.91\$ (91\% of students are encouraged).} \\
\texttt{    *   Consequently, \$P(\textbackslash{}neg E) = 1 {-} 0.91 = 0.09\$ (9\% are not encouraged).} \\
\texttt{    *   The problem states the probability of a high score given encouragement: \$P(S|E) = 0.52\$.} \\
\texttt{    *   The problem states the probability of a high score given no encouragement: \$P(S|\textbackslash{}neg E) = 0.72\$.} \\
\\ 
\texttt{    *Note: The intermediate variable "studying habit" acts as a mediator in the causal chain (\$E \textbackslash{}to \textbackslash{}text\{Habit\} \textbackslash{}to S\$), but since the problem provides the direct conditional probabilities of the outcome (\$S\$) given the treatment (\$E\$), we can bypass the specific distribution of the habit to calculate the overall outcome probability using the Law of Total Probability.*} \\
\\ 
\texttt{2.  **Calculate the Overall Probability of a High Score (\$P(S)\$):**} \\
\texttt{    Using the Law of Total Probability:} \\
\texttt{    \$\$P(S) = P(S|E) \textbackslash{}cdot P(E) + P(S|\textbackslash{}neg E) \textbackslash{}cdot P(\textbackslash{}neg E)\$\$} \\
\\ 
\texttt{    Substituting the values:} \\
\texttt{    \$\$P(S) = (0.52 \textbackslash{}times 0.91) + (0.72 \textbackslash{}times 0.09)\$\$} \\
\texttt{    \$\$P(S) = 0.4732 + 0.0648\$\$} \\
\texttt{    \$\$P(S) = 0.538\$\$} \\
\\ 
\texttt{3.  **Compare Probabilities:**} \\
\texttt{    *   Probability of High Score: \$0.538\$ (53.8\%)} \\
\texttt{    *   Probability of Low Score: \$1 {-} 0.538 = 0.462\$ (46.2\%)} \\
\\ 
\texttt{Since \$0.538 > 0.462\$, the probability of a high exam score is greater than the probability of a low exam score.} \\
\\ 
\texttt{**Conclusion:**} \\
\texttt{Despite the counter{-}intuitive result where the treatment (encouragement) lowers the conditional probability of success compared to the untreated group (a classic case of **negative confounding** or **selection bias** where the "unencouraged" group is inherently more capable), the sheer size of} \\
            }
        }
    }
    & & \\ \\

    \theutterance \stepcounter{utterance}  
    & & & \multicolumn{2}{p{0.3\linewidth}}{
        \cellcolor[rgb]{0.9,0.9,0.9}{
            \makecell[{{p{\linewidth}}}]{
                \texttt{\tiny{[GM$|$GM]}}
                \texttt{Yes} \\
            }
        }
    }
    & & \\ \\

    \theutterance \stepcounter{utterance}  
    & & & \multicolumn{2}{p{0.3\linewidth}}{
        \cellcolor[rgb]{0.9,0.9,0.9}{
            \makecell[{{p{\linewidth}}}]{
                \texttt{\tiny{[GM$|$GM]}}
                \texttt{yes} \\
            }
        }
    }
    & & \\ \\

    \theutterance \stepcounter{utterance}  
    & & & \multicolumn{2}{p{0.3\linewidth}}{
        \cellcolor[rgb]{0.9,0.9,0.9}{
            \makecell[{{p{\linewidth}}}]{
                \texttt{\tiny{[GM$|$GM]}}
                \texttt{game\_result = WIN} \\
            }
        }
    }
    & & \\ \\

\end{supertabular}
}

\end{document}
